% Lösungen Einheit 4
\einheitenkopf*{Einheit 4 von 5 · Verdopplungs- und Halbwertszeit}

\erg{36}{a) Bestand \quad b) Zeit \quad c) Bestand \quad d) Zeit}
\erg{37}{a) $180$, Verdopplungszeit \quad b) $700$, Halbwertszeit \quad c) $1{,}25\,\text{mg}$, Halbwertszeit \quad d) $1\,450$, Verdopplungszeit \quad e) $8$ Tage ($800\,\text{g}$ ist die Hälfte von $1\,600\,\text{g}$) \quad f) etwa $6$ Jahre (der doppelte Wert $600$ liegt zwischen $553$ und $625$, näher an $625$)}
\erg{38}{a) $3$ Stunden \quad b) Ich suche an der senkrechten Achse die Hälfte des Anfangswerts, also $160\,\text{mg}$, und gehe waagerecht bis zum Graphen. Von dort gehe ich senkrecht nach unten zur Zeitachse und lese $3$ Stunden ab.}
\erg{39}{a) nach $5$ Jahren ($46;\ 52{,}9;\ 60{,}8;\ 70{,}0;\ 80{,}5$) \quad b) $345;\ 396{,}8;\ 456{,}3;\ 524{,}7;\ 603{,}4$ – auch nach $5$ Jahren erstmals über $600$ \quad c) $t \approx 4{,}96 \approx 5{,}0$ Jahre \quad d) Anfangswert $1{,}6$ Tausend, doppelt: $3{,}2$ Tausend; waagerecht zum Graphen, senkrecht nach unten: etwa $8$ Jahre}
\erg{40}{a) Fehler: Emil halbiert die Zeit statt der Menge. Gesucht ist die Zeit, nach der die Hälfte von $500\,\text{mg}$, also $250\,\text{mg}$, übrig ist: $8$ Stunden. \quad b) $18$ Jahre ($1\,200\,\text{g}$ ist die Hälfte von $2\,400\,\text{g}$)}
\erg{41}{a) drei Verdopplungen: $2 \cdot 2 \cdot 2 = 8$, also achtmal so groß \quad b) In jeder gleich langen Zeitspanne wird mit demselben Faktor malgenommen, egal wie groß der Bestand gerade ist; deshalb dauert jede Verdopplung gleich lang. \quad c) Nein: Derselbe Betrag ist bei einem größeren Bestand ein kleinerer Anteil; jede weitere Verdopplung dauert länger.}
\erg{42}{a) $4\,\text{kg}$; $2\,\text{kg}$; $1\,\text{kg}$ \quad b) Nein: Nach $24$ Tagen sind es $2\,\text{kg}$, nach $36$ Tagen $1\,\text{kg}$; nach $28$ Tagen liegt die Menge dazwischen (etwa $1{,}59\,\text{kg}$), also über $1\,\text{kg}$. \quad c) nach mehr als $36$ Tagen}
